\documentclass[11pt]{article}
% jugeo-common.tex — shared preamble for all Judgment Geometry papers
% Include via: % jugeo-common.tex — shared preamble for all Judgment Geometry papers
% Include via: % jugeo-common.tex — shared preamble for all Judgment Geometry papers
% Include via: \input{jugeo-common}

% ── Packages ────────────────────────────────────────────────────────────
\usepackage{amsmath,amssymb,amsthm,mathtools}
\usepackage{tikz-cd}
\usepackage{listings}
\usepackage{xcolor}
\usepackage{xspace}
\usepackage{booktabs}
\usepackage{multirow}
\usepackage{enumitem}
\usepackage[expansion=false]{microtype}
\usepackage{hyperref}
\usepackage{cleveref}
% \usepackage{thmtools}  % disabled: not available in this TeX installation
\usepackage{float}
\usepackage{caption}
\usepackage{subcaption}
\usepackage{tabularx}
\usepackage{stmaryrd}       % \llbracket, \rrbracket
\usepackage{bbm}            % \mathbbm{1}

% ── Page geometry ───────────────────────────────────────────────────────
\usepackage[margin=1in]{geometry}

% ── Theorem environments ────────────────────────────────────────────────
\theoremstyle{plain}
\newtheorem{theorem}{Theorem}[section]
\newtheorem{lemma}[theorem]{Lemma}
\newtheorem{proposition}[theorem]{Proposition}
\newtheorem{corollary}[theorem]{Corollary}

\theoremstyle{definition}
\newtheorem{definition}[theorem]{Definition}
\newtheorem{example}[theorem]{Example}
\newtheorem{construction}[theorem]{Construction}

\theoremstyle{remark}
\newtheorem{remark}[theorem]{Remark}
\newtheorem{notation}[theorem]{Notation}

% ── Python listings style ──────────────────────────────────────────────
\definecolor{jugeoBlue}{HTML}{2563EB}
\definecolor{jugeoGreen}{HTML}{059669}
\definecolor{jugeoRed}{HTML}{DC2626}
\definecolor{jugeoPurple}{HTML}{7C3AED}
\definecolor{jugeoGray}{HTML}{6B7280}
\definecolor{jugeoBg}{HTML}{F8FAFC}

\lstdefinestyle{jugeo-python}{
  language=Python,
  basicstyle=\ttfamily\small,
  keywordstyle=\color{jugeoBlue}\bfseries,
  stringstyle=\color{jugeoGreen},
  commentstyle=\color{jugeoGray}\itshape,
  numberstyle=\tiny\color{jugeoGray},
  backgroundcolor=\color{jugeoBg},
  numbers=left,
  numbersep=6pt,
  frame=single,
  framerule=0.4pt,
  rulecolor=\color{jugeoGray!40},
  breaklines=true,
  breakatwhitespace=true,
  showstringspaces=false,
  tabsize=4,
  xleftmargin=1.5em,
  framexleftmargin=1.5em,
  morekeywords={dataclass,frozen,field,Enum,IntEnum,auto,Any,
                Mapping,Sequence,Optional,match,case},
  literate={→}{{$\to$}}1 {←}{{$\leftarrow$}}1
           {≤}{{$\leq$}}1 {≥}{{$\geq$}}1 {≠}{{$\neq$}}1
           {∈}{{$\in$}}1 {∀}{{$\forall$}}1 {∃}{{$\exists$}}1
           {⊕}{{$\oplus$}}1 {⊖}{{$\ominus$}}1,
  escapeinside={(*@}{@*)},
}
\lstset{style=jugeo-python}

\lstdefinestyle{jugeo-output}{
  basicstyle=\ttfamily\small,
  backgroundcolor=\color{jugeoBg},
  frame=single,
  framerule=0.4pt,
  rulecolor=\color{jugeoGray!40},
  breaklines=true,
  showstringspaces=false,
  xleftmargin=1.5em,
  framexleftmargin=0.5em,
  numbers=none,
}

% ── JuGeo-specific macros ──────────────────────────────────────────────

% Judgment 8-tuple
\newcommand{\J}{\mathcal{J}}
\newcommand{\Jud}[8]{(#1,\,#2,\,#3,\,#4,\,#5,\,#6,\,#7,\,#8)}
\newcommand{\JudShort}[1]{\J_{#1}}

% Components
\newcommand{\coord}{\mathsf{c}}            % coordinate
\newcommand{\prop}{\varphi}                 % proposition
\newcommand{\carrier}{\mathcal{A}}          % carrier type
\newcommand{\evid}{\mathcal{E}}             % evidence bundle
\newcommand{\oblig}{\mathcal{O}}            % obligations
\newcommand{\obstr}{\mathcal{B}}            % obstructions
\newcommand{\trust}{\mathcal{T}}            % trust annotation
\newcommand{\prov}{\Pi}                     % provenance

% Trust algebra
\newcommand{\Talg}{\mathfrak{T}}            % trust ordered algebra
\newcommand{\Eadm}{\mathcal{E}_{\mathrm{adm}}}
\newcommand{\tleq}{\preceq}                 % trust ordering
\newcommand{\tjoin}{\oplus}                 % conservative join
\newcommand{\tatten}{\ominus}               % attenuation
\newcommand{\tpromote}{\uparrow_\pi}        % promotion
\newcommand{\tdemote}{\downarrow_\chi}       % demotion

% Trust tiers
\newcommand{\Tcontra}{\ensuremath{\mathsf{CONTRADICTED}}}
\newcommand{\Tunver}{\ensuremath{\mathsf{UNVERIFIED}}}
\newcommand{\Tcopilot}{\ensuremath{\mathsf{COPILOT}}}
\newcommand{\Toracle}{\ensuremath{\mathsf{ORACLE}}}
\newcommand{\Truntime}{\ensuremath{\mathsf{RUNTIME}}}
\newcommand{\Tsolver}{\ensuremath{\mathsf{SOLVER}}}
\newcommand{\Tproof}{\ensuremath{\mathsf{PROOF}}}

% Site and geometry
\newcommand{\Site}{\mathcal{S}}             % semantic site
\newcommand{\Cov}{\mathrm{Cov}}            % covering family
\newcommand{\Ob}{\mathrm{Ob}}              % objects
\newcommand{\Mor}{\mathrm{Mor}}            % morphisms
\newcommand{\res}{\mathrm{res}}            % restriction
\newcommand{\Cech}{\check{C}}              % Čech complex
\newcommand{\Hcoh}[1]{H^{#1}}             % cohomology group

% Descent
\newcommand{\descent}{\mathrm{desc}}
\newcommand{\glue}{\mathrm{glue}}
\newcommand{\overlap}[2]{#1 \cap #2}

% Semantic moves
\newcommand{\VERIFY}{\textsc{Verify}}
\newcommand{\CONSTRUCT}{\textsc{Construct}}
\newcommand{\NORMALIZE}{\textsc{Normalize}}
\newcommand{\GLUE}{\textsc{Glue}}
\newcommand{\RESTRICT}{\textsc{Restrict}}
\newcommand{\TRANSPORT}{\textsc{Transport}}
\newcommand{\REPAIR}{\textsc{Repair}}
\newcommand{\PROMOTE}{\textsc{Promote}}
\newcommand{\CHALLENGE}{\textsc{Challenge}}
\newcommand{\ESCALATE}{\textsc{Escalate}}

% Fragments
\newcommand{\QFLIA}{\texttt{QF\_LIA}}
\newcommand{\QFLRA}{\texttt{QF\_LRA}}
\newcommand{\QFBV}{\texttt{QF\_BV}}
\newcommand{\QFUF}{\texttt{QF\_UF}}

% Misc
\newcommand{\Python}{\textsc{Python}}
\newcommand{\JuGeo}{\textsc{JuGeo}}
\newcommand{\LEAN}{\textsc{Lean}\,4}
\newcommand{\Fstar}{F$^{\star}$}
\newcommand{\Coq}{\textsc{Coq}}
\newcommand{\Dafny}{\textsc{Dafny}}

% Common shortcuts
\newcommand{\ie}{i.e.\@\xspace}
\newcommand{\eg}{e.g.\@\xspace}
\newcommand{\cf}{cf.\@\xspace}
\newcommand{\wrt}{w.r.t.\@\xspace}

% Evaluation shortcuts
\newcommand{\numcases}{300}
\newcommand{\numspec}{100}
\newcommand{\numeq}{100}
\newcommand{\numbug}{100}
\newcommand{\medtime}{6.5\,ms}
\newcommand{\totaltime}{1.949\,s}


% ── Packages ────────────────────────────────────────────────────────────
\usepackage{amsmath,amssymb,amsthm,mathtools}
\usepackage{tikz-cd}
\usepackage{listings}
\usepackage{xcolor}
\usepackage{xspace}
\usepackage{booktabs}
\usepackage{multirow}
\usepackage{enumitem}
\usepackage[expansion=false]{microtype}
\usepackage{hyperref}
\usepackage{cleveref}
% \usepackage{thmtools}  % disabled: not available in this TeX installation
\usepackage{float}
\usepackage{caption}
\usepackage{subcaption}
\usepackage{tabularx}
\usepackage{stmaryrd}       % \llbracket, \rrbracket
\usepackage{bbm}            % \mathbbm{1}

% ── Page geometry ───────────────────────────────────────────────────────
\usepackage[margin=1in]{geometry}

% ── Theorem environments ────────────────────────────────────────────────
\theoremstyle{plain}
\newtheorem{theorem}{Theorem}[section]
\newtheorem{lemma}[theorem]{Lemma}
\newtheorem{proposition}[theorem]{Proposition}
\newtheorem{corollary}[theorem]{Corollary}

\theoremstyle{definition}
\newtheorem{definition}[theorem]{Definition}
\newtheorem{example}[theorem]{Example}
\newtheorem{construction}[theorem]{Construction}

\theoremstyle{remark}
\newtheorem{remark}[theorem]{Remark}
\newtheorem{notation}[theorem]{Notation}

% ── Python listings style ──────────────────────────────────────────────
\definecolor{jugeoBlue}{HTML}{2563EB}
\definecolor{jugeoGreen}{HTML}{059669}
\definecolor{jugeoRed}{HTML}{DC2626}
\definecolor{jugeoPurple}{HTML}{7C3AED}
\definecolor{jugeoGray}{HTML}{6B7280}
\definecolor{jugeoBg}{HTML}{F8FAFC}

\lstdefinestyle{jugeo-python}{
  language=Python,
  basicstyle=\ttfamily\small,
  keywordstyle=\color{jugeoBlue}\bfseries,
  stringstyle=\color{jugeoGreen},
  commentstyle=\color{jugeoGray}\itshape,
  numberstyle=\tiny\color{jugeoGray},
  backgroundcolor=\color{jugeoBg},
  numbers=left,
  numbersep=6pt,
  frame=single,
  framerule=0.4pt,
  rulecolor=\color{jugeoGray!40},
  breaklines=true,
  breakatwhitespace=true,
  showstringspaces=false,
  tabsize=4,
  xleftmargin=1.5em,
  framexleftmargin=1.5em,
  morekeywords={dataclass,frozen,field,Enum,IntEnum,auto,Any,
                Mapping,Sequence,Optional,match,case},
  literate={→}{{$\to$}}1 {←}{{$\leftarrow$}}1
           {≤}{{$\leq$}}1 {≥}{{$\geq$}}1 {≠}{{$\neq$}}1
           {∈}{{$\in$}}1 {∀}{{$\forall$}}1 {∃}{{$\exists$}}1
           {⊕}{{$\oplus$}}1 {⊖}{{$\ominus$}}1,
  escapeinside={(*@}{@*)},
}
\lstset{style=jugeo-python}

\lstdefinestyle{jugeo-output}{
  basicstyle=\ttfamily\small,
  backgroundcolor=\color{jugeoBg},
  frame=single,
  framerule=0.4pt,
  rulecolor=\color{jugeoGray!40},
  breaklines=true,
  showstringspaces=false,
  xleftmargin=1.5em,
  framexleftmargin=0.5em,
  numbers=none,
}

% ── JuGeo-specific macros ──────────────────────────────────────────────

% Judgment 8-tuple
\newcommand{\J}{\mathcal{J}}
\newcommand{\Jud}[8]{(#1,\,#2,\,#3,\,#4,\,#5,\,#6,\,#7,\,#8)}
\newcommand{\JudShort}[1]{\J_{#1}}

% Components
\newcommand{\coord}{\mathsf{c}}            % coordinate
\newcommand{\prop}{\varphi}                 % proposition
\newcommand{\carrier}{\mathcal{A}}          % carrier type
\newcommand{\evid}{\mathcal{E}}             % evidence bundle
\newcommand{\oblig}{\mathcal{O}}            % obligations
\newcommand{\obstr}{\mathcal{B}}            % obstructions
\newcommand{\trust}{\mathcal{T}}            % trust annotation
\newcommand{\prov}{\Pi}                     % provenance

% Trust algebra
\newcommand{\Talg}{\mathfrak{T}}            % trust ordered algebra
\newcommand{\Eadm}{\mathcal{E}_{\mathrm{adm}}}
\newcommand{\tleq}{\preceq}                 % trust ordering
\newcommand{\tjoin}{\oplus}                 % conservative join
\newcommand{\tatten}{\ominus}               % attenuation
\newcommand{\tpromote}{\uparrow_\pi}        % promotion
\newcommand{\tdemote}{\downarrow_\chi}       % demotion

% Trust tiers
\newcommand{\Tcontra}{\ensuremath{\mathsf{CONTRADICTED}}}
\newcommand{\Tunver}{\ensuremath{\mathsf{UNVERIFIED}}}
\newcommand{\Tcopilot}{\ensuremath{\mathsf{COPILOT}}}
\newcommand{\Toracle}{\ensuremath{\mathsf{ORACLE}}}
\newcommand{\Truntime}{\ensuremath{\mathsf{RUNTIME}}}
\newcommand{\Tsolver}{\ensuremath{\mathsf{SOLVER}}}
\newcommand{\Tproof}{\ensuremath{\mathsf{PROOF}}}

% Site and geometry
\newcommand{\Site}{\mathcal{S}}             % semantic site
\newcommand{\Cov}{\mathrm{Cov}}            % covering family
\newcommand{\Ob}{\mathrm{Ob}}              % objects
\newcommand{\Mor}{\mathrm{Mor}}            % morphisms
\newcommand{\res}{\mathrm{res}}            % restriction
\newcommand{\Cech}{\check{C}}              % Čech complex
\newcommand{\Hcoh}[1]{H^{#1}}             % cohomology group

% Descent
\newcommand{\descent}{\mathrm{desc}}
\newcommand{\glue}{\mathrm{glue}}
\newcommand{\overlap}[2]{#1 \cap #2}

% Semantic moves
\newcommand{\VERIFY}{\textsc{Verify}}
\newcommand{\CONSTRUCT}{\textsc{Construct}}
\newcommand{\NORMALIZE}{\textsc{Normalize}}
\newcommand{\GLUE}{\textsc{Glue}}
\newcommand{\RESTRICT}{\textsc{Restrict}}
\newcommand{\TRANSPORT}{\textsc{Transport}}
\newcommand{\REPAIR}{\textsc{Repair}}
\newcommand{\PROMOTE}{\textsc{Promote}}
\newcommand{\CHALLENGE}{\textsc{Challenge}}
\newcommand{\ESCALATE}{\textsc{Escalate}}

% Fragments
\newcommand{\QFLIA}{\texttt{QF\_LIA}}
\newcommand{\QFLRA}{\texttt{QF\_LRA}}
\newcommand{\QFBV}{\texttt{QF\_BV}}
\newcommand{\QFUF}{\texttt{QF\_UF}}

% Misc
\newcommand{\Python}{\textsc{Python}}
\newcommand{\JuGeo}{\textsc{JuGeo}}
\newcommand{\LEAN}{\textsc{Lean}\,4}
\newcommand{\Fstar}{F$^{\star}$}
\newcommand{\Coq}{\textsc{Coq}}
\newcommand{\Dafny}{\textsc{Dafny}}

% Common shortcuts
\newcommand{\ie}{i.e.\@\xspace}
\newcommand{\eg}{e.g.\@\xspace}
\newcommand{\cf}{cf.\@\xspace}
\newcommand{\wrt}{w.r.t.\@\xspace}

% Evaluation shortcuts
\newcommand{\numcases}{300}
\newcommand{\numspec}{100}
\newcommand{\numeq}{100}
\newcommand{\numbug}{100}
\newcommand{\medtime}{6.5\,ms}
\newcommand{\totaltime}{1.949\,s}


% ── Packages ────────────────────────────────────────────────────────────
\usepackage{amsmath,amssymb,amsthm,mathtools}
\usepackage{tikz-cd}
\usepackage{listings}
\usepackage{xcolor}
\usepackage{xspace}
\usepackage{booktabs}
\usepackage{multirow}
\usepackage{enumitem}
\usepackage[expansion=false]{microtype}
\usepackage{hyperref}
\usepackage{cleveref}
% \usepackage{thmtools}  % disabled: not available in this TeX installation
\usepackage{float}
\usepackage{caption}
\usepackage{subcaption}
\usepackage{tabularx}
\usepackage{stmaryrd}       % \llbracket, \rrbracket
\usepackage{bbm}            % \mathbbm{1}

% ── Page geometry ───────────────────────────────────────────────────────
\usepackage[margin=1in]{geometry}

% ── Theorem environments ────────────────────────────────────────────────
\theoremstyle{plain}
\newtheorem{theorem}{Theorem}[section]
\newtheorem{lemma}[theorem]{Lemma}
\newtheorem{proposition}[theorem]{Proposition}
\newtheorem{corollary}[theorem]{Corollary}

\theoremstyle{definition}
\newtheorem{definition}[theorem]{Definition}
\newtheorem{example}[theorem]{Example}
\newtheorem{construction}[theorem]{Construction}

\theoremstyle{remark}
\newtheorem{remark}[theorem]{Remark}
\newtheorem{notation}[theorem]{Notation}

% ── Python listings style ──────────────────────────────────────────────
\definecolor{jugeoBlue}{HTML}{2563EB}
\definecolor{jugeoGreen}{HTML}{059669}
\definecolor{jugeoRed}{HTML}{DC2626}
\definecolor{jugeoPurple}{HTML}{7C3AED}
\definecolor{jugeoGray}{HTML}{6B7280}
\definecolor{jugeoBg}{HTML}{F8FAFC}

\lstdefinestyle{jugeo-python}{
  language=Python,
  basicstyle=\ttfamily\small,
  keywordstyle=\color{jugeoBlue}\bfseries,
  stringstyle=\color{jugeoGreen},
  commentstyle=\color{jugeoGray}\itshape,
  numberstyle=\tiny\color{jugeoGray},
  backgroundcolor=\color{jugeoBg},
  numbers=left,
  numbersep=6pt,
  frame=single,
  framerule=0.4pt,
  rulecolor=\color{jugeoGray!40},
  breaklines=true,
  breakatwhitespace=true,
  showstringspaces=false,
  tabsize=4,
  xleftmargin=1.5em,
  framexleftmargin=1.5em,
  morekeywords={dataclass,frozen,field,Enum,IntEnum,auto,Any,
                Mapping,Sequence,Optional,match,case},
  literate={→}{{$\to$}}1 {←}{{$\leftarrow$}}1
           {≤}{{$\leq$}}1 {≥}{{$\geq$}}1 {≠}{{$\neq$}}1
           {∈}{{$\in$}}1 {∀}{{$\forall$}}1 {∃}{{$\exists$}}1
           {⊕}{{$\oplus$}}1 {⊖}{{$\ominus$}}1,
  escapeinside={(*@}{@*)},
}
\lstset{style=jugeo-python}

\lstdefinestyle{jugeo-output}{
  basicstyle=\ttfamily\small,
  backgroundcolor=\color{jugeoBg},
  frame=single,
  framerule=0.4pt,
  rulecolor=\color{jugeoGray!40},
  breaklines=true,
  showstringspaces=false,
  xleftmargin=1.5em,
  framexleftmargin=0.5em,
  numbers=none,
}

% ── JuGeo-specific macros ──────────────────────────────────────────────

% Judgment 8-tuple
\newcommand{\J}{\mathcal{J}}
\newcommand{\Jud}[8]{(#1,\,#2,\,#3,\,#4,\,#5,\,#6,\,#7,\,#8)}
\newcommand{\JudShort}[1]{\J_{#1}}

% Components
\newcommand{\coord}{\mathsf{c}}            % coordinate
\newcommand{\prop}{\varphi}                 % proposition
\newcommand{\carrier}{\mathcal{A}}          % carrier type
\newcommand{\evid}{\mathcal{E}}             % evidence bundle
\newcommand{\oblig}{\mathcal{O}}            % obligations
\newcommand{\obstr}{\mathcal{B}}            % obstructions
\newcommand{\trust}{\mathcal{T}}            % trust annotation
\newcommand{\prov}{\Pi}                     % provenance

% Trust algebra
\newcommand{\Talg}{\mathfrak{T}}            % trust ordered algebra
\newcommand{\Eadm}{\mathcal{E}_{\mathrm{adm}}}
\newcommand{\tleq}{\preceq}                 % trust ordering
\newcommand{\tjoin}{\oplus}                 % conservative join
\newcommand{\tatten}{\ominus}               % attenuation
\newcommand{\tpromote}{\uparrow_\pi}        % promotion
\newcommand{\tdemote}{\downarrow_\chi}       % demotion

% Trust tiers
\newcommand{\Tcontra}{\ensuremath{\mathsf{CONTRADICTED}}}
\newcommand{\Tunver}{\ensuremath{\mathsf{UNVERIFIED}}}
\newcommand{\Tcopilot}{\ensuremath{\mathsf{COPILOT}}}
\newcommand{\Toracle}{\ensuremath{\mathsf{ORACLE}}}
\newcommand{\Truntime}{\ensuremath{\mathsf{RUNTIME}}}
\newcommand{\Tsolver}{\ensuremath{\mathsf{SOLVER}}}
\newcommand{\Tproof}{\ensuremath{\mathsf{PROOF}}}

% Site and geometry
\newcommand{\Site}{\mathcal{S}}             % semantic site
\newcommand{\Cov}{\mathrm{Cov}}            % covering family
\newcommand{\Ob}{\mathrm{Ob}}              % objects
\newcommand{\Mor}{\mathrm{Mor}}            % morphisms
\newcommand{\res}{\mathrm{res}}            % restriction
\newcommand{\Cech}{\check{C}}              % Čech complex
\newcommand{\Hcoh}[1]{H^{#1}}             % cohomology group

% Descent
\newcommand{\descent}{\mathrm{desc}}
\newcommand{\glue}{\mathrm{glue}}
\newcommand{\overlap}[2]{#1 \cap #2}

% Semantic moves
\newcommand{\VERIFY}{\textsc{Verify}}
\newcommand{\CONSTRUCT}{\textsc{Construct}}
\newcommand{\NORMALIZE}{\textsc{Normalize}}
\newcommand{\GLUE}{\textsc{Glue}}
\newcommand{\RESTRICT}{\textsc{Restrict}}
\newcommand{\TRANSPORT}{\textsc{Transport}}
\newcommand{\REPAIR}{\textsc{Repair}}
\newcommand{\PROMOTE}{\textsc{Promote}}
\newcommand{\CHALLENGE}{\textsc{Challenge}}
\newcommand{\ESCALATE}{\textsc{Escalate}}

% Fragments
\newcommand{\QFLIA}{\texttt{QF\_LIA}}
\newcommand{\QFLRA}{\texttt{QF\_LRA}}
\newcommand{\QFBV}{\texttt{QF\_BV}}
\newcommand{\QFUF}{\texttt{QF\_UF}}

% Misc
\newcommand{\Python}{\textsc{Python}}
\newcommand{\JuGeo}{\textsc{JuGeo}}
\newcommand{\LEAN}{\textsc{Lean}\,4}
\newcommand{\Fstar}{F$^{\star}$}
\newcommand{\Coq}{\textsc{Coq}}
\newcommand{\Dafny}{\textsc{Dafny}}

% Common shortcuts
\newcommand{\ie}{i.e.\@\xspace}
\newcommand{\eg}{e.g.\@\xspace}
\newcommand{\cf}{cf.\@\xspace}
\newcommand{\wrt}{w.r.t.\@\xspace}

% Evaluation shortcuts
\newcommand{\numcases}{300}
\newcommand{\numspec}{100}
\newcommand{\numeq}{100}
\newcommand{\numbug}{100}
\newcommand{\medtime}{6.5\,ms}
\newcommand{\totaltime}{1.949\,s}

% experiment-data.tex — AUTO-GENERATED from real jugeo CLI runs
% DO NOT EDIT — regenerate with: python3 experiments/run_universal_metrics.py
% Generated from 30 programs across 6 domains

\newcommand{\expTotalPrograms}{30}
\newcommand{\expVerified}{30}
\newcommand{\expAccuracy}{100.0\%}
\newcommand{\expCoordsMin}{2}
\newcommand{\expCoordsMax}{10}
\newcommand{\expCoordsMean}{5.2}
\newcommand{\expCoordsMedian}{5.5}
\newcommand{\expPropsMin}{4}
\newcommand{\expPropsMax}{20}
\newcommand{\expPropsMean}{10.4}
\newcommand{\expPropsSum}{311}
\newcommand{\expPropsOkSum}{310}
\newcommand{\expTimeMin}{3.506\,s}
\newcommand{\expTimeMax}{6.714\,s}
\newcommand{\expTimeMean}{4.64\,s}
\newcommand{\expTimeMedian}{4.508\,s}
\newcommand{\expTimeTotal}{139.204\,s}
\newcommand{\expObstructionTotal}{0}
\newcommand{\expHOneZero}{30}
\newcommand{\expDomainCount}{6}
\newcommand{\expTrustCOPILOTSUGGESTED}{30}
\newcommand{\expDomainDsCount}{5}
\newcommand{\expDomainDsVerified}{5}
\newcommand{\expDomainDsAvgCoords}{7.6}
\newcommand{\expDomainDsAvgProps}{15.2}
\newcommand{\expDomainMathCount}{5}
\newcommand{\expDomainMathVerified}{5}
\newcommand{\expDomainMathAvgCoords}{4.8}
\newcommand{\expDomainMathAvgProps}{9.4}
\newcommand{\expDomainSmCount}{5}
\newcommand{\expDomainSmVerified}{5}
\newcommand{\expDomainSmAvgCoords}{7.6}
\newcommand{\expDomainSmAvgProps}{15.2}
\newcommand{\expDomainSortCount}{5}
\newcommand{\expDomainSortVerified}{5}
\newcommand{\expDomainSortAvgCoords}{2.4}
\newcommand{\expDomainSortAvgProps}{4.8}
\newcommand{\expDomainStrCount}{5}
\newcommand{\expDomainStrVerified}{5}
\newcommand{\expDomainStrAvgCoords}{3.2}
\newcommand{\expDomainStrAvgProps}{6.4}
\newcommand{\expDomainWebCount}{5}
\newcommand{\expDomainWebVerified}{5}
\newcommand{\expDomainWebAvgCoords}{5.6}
\newcommand{\expDomainWebAvgProps}{11.2}

% subsystem-data.tex — AUTO-GENERATED from real jugeo subsystem experiments
% DO NOT EDIT — regenerate with: python3 experiments/run_subsystem_experiments.py

\newcommand{\subCliTotal}{10}
\newcommand{\subCliVerified}{9}
\newcommand{\subCliAccuracy}{90.0\%}
\newcommand{\subCliMeanTime}{4.132\,s}
\newcommand{\subCliMinTime}{3.472\,s}
\newcommand{\subCliMaxTime}{5.372\,s}
\newcommand{\subCliTotalCoords}{54}
\newcommand{\subCliTotalProps}{107}
\newcommand{\subCliTotalPropsOk}{106}
\newcommand{\subSiteCalls}{90}
\newcommand{\subSiteSuccessful}{69}
\newcommand{\subSiteSuccessRate}{76.7\%}
\newcommand{\subSiteMeanTime}{0.0001\,s}
\newcommand{\subSiteTotalTime}{0.005\,s}
\newcommand{\subSpecificationsatisfactionSmallTime}{0.0003\,s}
\newcommand{\subSpecificationsatisfactionMediumTime}{0.0001\,s}
\newcommand{\subSpecificationsatisfactionLargeTime}{0.0001\,s}
\newcommand{\subInhabitantfleetSmallTime}{0.0\,s}
\newcommand{\subInhabitantfleetMediumTime}{0.0\,s}
\newcommand{\subInhabitantfleetLargeTime}{0.0\,s}
\newcommand{\subTheoremecologySmallTime}{0.0\,s}
\newcommand{\subTheoremecologyMediumTime}{0.0\,s}
\newcommand{\subTheoremecologyLargeTime}{0.0\,s}
\newcommand{\subStatespaceexplorationSmallTime}{0.0\,s}
\newcommand{\subStatespaceexplorationMediumTime}{0.0\,s}
\newcommand{\subStatespaceexplorationLargeTime}{0.0\,s}
\newcommand{\subRepairsemanticsSmallTime}{0.0\,s}
\newcommand{\subRepairsemanticsMediumTime}{0.0\,s}
\newcommand{\subRepairsemanticsLargeTime}{0.0\,s}
\newcommand{\subReplaygluingSmallTime}{0.0\,s}
\newcommand{\subReplaygluingMediumTime}{0.0\,s}
\newcommand{\subReplaygluingLargeTime}{0.0\,s}
\newcommand{\subSemanticfuturesSmallTime}{0.0\,s}
\newcommand{\subSemanticfuturesMediumTime}{0.0\,s}
\newcommand{\subSemanticfuturesLargeTime}{0.0\,s}
\newcommand{\subHypercovertreatySmallTime}{0.0\,s}
\newcommand{\subHypercovertreatyMediumTime}{0.0\,s}
\newcommand{\subHypercovertreatyLargeTime}{0.0\,s}
\newcommand{\subEncodeforsolverSmallTime}{0.0026\,s}
\newcommand{\subEncodeforsolverMediumTime}{0.0002\,s}
\newcommand{\subEncodeforsolverLargeTime}{0.0001\,s}
\newcommand{\subInterfaceroutingSmallTime}{0.0\,s}
\newcommand{\subInterfaceroutingMediumTime}{0.0\,s}
\newcommand{\subInterfaceroutingLargeTime}{0.0\,s}
\newcommand{\subOrchestrateverificationSmallTime}{0.0002\,s}
\newcommand{\subOrchestrateverificationMediumTime}{0.0\,s}
\newcommand{\subOrchestrateverificationLargeTime}{0.0\,s}
\newcommand{\subRunfulldescentSmallTime}{0.0\,s}
\newcommand{\subRunfulldescentMediumTime}{0.0\,s}
\newcommand{\subRunfulldescentLargeTime}{0.0\,s}
\newcommand{\subKernellifecycleSmallTime}{0.0\,s}
\newcommand{\subKernellifecycleMediumTime}{0.0\,s}
\newcommand{\subKernellifecycleLargeTime}{0.0\,s}
\newcommand{\subDiscoverypipelineSmallTime}{0.0\,s}
\newcommand{\subDiscoverypipelineMediumTime}{0.0\,s}
\newcommand{\subDiscoverypipelineLargeTime}{0.0\,s}
\newcommand{\subAnalogytransportSmallTime}{0.0\,s}
\newcommand{\subAnalogytransportMediumTime}{0.0\,s}
\newcommand{\subAnalogytransportLargeTime}{0.0\,s}
\newcommand{\subFormalcoresiteSmallTime}{0.0001\,s}
\newcommand{\subFormalcoresiteMediumTime}{0.0\,s}
\newcommand{\subFormalcoresiteLargeTime}{0.0\,s}
\newcommand{\subProblematlasSmallTime}{0.0\,s}
\newcommand{\subProblematlasMediumTime}{0.0\,s}
\newcommand{\subProblematlasLargeTime}{0.0\,s}
\newcommand{\subPublicalignmentSmallTime}{0.0\,s}
\newcommand{\subPublicalignmentMediumTime}{0.0\,s}
\newcommand{\subPublicalignmentLargeTime}{0.0\,s}
\newcommand{\subRegimebootstrappingSmallTime}{0.0\,s}
\newcommand{\subRegimebootstrappingMediumTime}{0.0\,s}
\newcommand{\subRegimebootstrappingLargeTime}{0.0\,s}
\newcommand{\subRelationalrefinementSmallTime}{0.0\,s}
\newcommand{\subRelationalrefinementMediumTime}{0.0\,s}
\newcommand{\subRelationalrefinementLargeTime}{0.0\,s}
\newcommand{\subSemanticclosureSmallTime}{0.0\,s}
\newcommand{\subSemanticclosureMediumTime}{0.0\,s}
\newcommand{\subSemanticclosureLargeTime}{0.0\,s}
\newcommand{\subTheoremeconomicsSmallTime}{0.0001\,s}
\newcommand{\subTheoremeconomicsMediumTime}{0.0\,s}
\newcommand{\subTheoremeconomicsLargeTime}{0.0\,s}
\newcommand{\subBenchmarksuiteSmallTime}{0.0\,s}
\newcommand{\subBenchmarksuiteMediumTime}{0.0\,s}
\newcommand{\subBenchmarksuiteLargeTime}{0.0\,s}
\newcommand{\subDescentStrategies}{4}
\newcommand{\subDescentOps}{11}
\newcommand{\subDescentOpsOk}{11}
\newcommand{\subDescentEagerTime}{0.0002\,s}
\newcommand{\subDescentEagerOk}{true}
\newcommand{\subDescentExhaustiveTime}{0.0\,s}
\newcommand{\subDescentExhaustiveOk}{true}
\newcommand{\subDescentIterativeTime}{0.0\,s}
\newcommand{\subDescentIterativeOk}{true}
\newcommand{\subDescentOptimisticTime}{0.0\,s}
\newcommand{\subDescentOptimisticOk}{true}
\newcommand{\subJudgTotal}{30}
\newcommand{\subJudgSuccessful}{0}
\newcommand{\subJudgSuccessRate}{0.0\%}
\newcommand{\subJudgMeanTimeUs}{7.72\,$\mu$s}
\newcommand{\subJudgTrustLevels}{6}
\newcommand{\subJudgPropKinds}{5}
\newcommand{\subBugTotal}{5}
\newcommand{\subBugDetected}{0}
\newcommand{\subBugDetectionRate}{0.0\%}
\newcommand{\subBugFalsePositiveRate}{0.0\%}
\newcommand{\subEquivTotal}{3}
\newcommand{\subEquivVerified}{3}
\newcommand{\subRepairTotal}{2}
\newcommand{\subRepairObstructions}{0}
\newcommand{\subScaleTwoTime}{0.0001\,s}
\newcommand{\subScaleFourTime}{0.0\,s}
\newcommand{\subScaleEightTime}{0.0\,s}
\newcommand{\subScaleTwelveTime}{0.0\,s}
\newcommand{\subScaleSixteenTime}{0.0\,s}
\newcommand{\subScaleTwentyTime}{0.0\,s}
\newcommand{\subTotalExperimentTime}{95.6\,s}

% comprehensive-data.tex — AUTO-GENERATED from real jugeo experiments
% DO NOT EDIT — regenerate with: python3 experiments/run_comprehensive_experiments.py
% Generated: 2026-03-23 09:19:06

% --- Size scaling metrics ---
\newcommand{\compTotalPrograms}{18}
\newcommand{\compTotalVerified}{18}
\newcommand{\compOverallAccuracy}{100.0\%}
\newcommand{\compTinyCount}{5}
\newcommand{\compTinyVerified}{5}
\newcommand{\compTinyMeanCoords}{3}
\newcommand{\compTinyMeanProps}{6}
\newcommand{\compTinyMeanTime}{4.95\,s}
\newcommand{\compTinyMeanLines}{5}
\newcommand{\compTinyMinTime}{4.45\,s}
\newcommand{\compTinyMaxTime}{5.51\,s}
\newcommand{\compSmallCount}{5}
\newcommand{\compSmallVerified}{5}
\newcommand{\compSmallMeanCoords}{3.6}
\newcommand{\compSmallMeanProps}{7.2}
\newcommand{\compSmallMeanTime}{4.85\,s}
\newcommand{\compSmallMeanLines}{16.0}
\newcommand{\compSmallMinTime}{4.30\,s}
\newcommand{\compSmallMaxTime}{5.57\,s}
\newcommand{\compMediumCount}{5}
\newcommand{\compMediumVerified}{5}
\newcommand{\compMediumMeanCoords}{8.6}
\newcommand{\compMediumMeanProps}{17.2}
\newcommand{\compMediumMeanTime}{4.47\,s}
\newcommand{\compMediumMeanLines}{45.0}
\newcommand{\compMediumMinTime}{4.26\,s}
\newcommand{\compMediumMaxTime}{4.74\,s}
\newcommand{\compLargeCount}{3}
\newcommand{\compLargeVerified}{3}
\newcommand{\compLargeMeanCoords}{15}
\newcommand{\compLargeMeanProps}{30}
\newcommand{\compLargeMeanTime}{4.31\,s}
\newcommand{\compLargeMeanLines}{79.0}
\newcommand{\compLargeMinTime}{4.27\,s}
\newcommand{\compLargeMaxTime}{4.38\,s}
% --- Aggregate timing ---
\newcommand{\compTimeMean}{4.68\,s}
\newcommand{\compTimeMedian}{4.50\,s}
\newcommand{\compTimeMin}{4.265\,s}
\newcommand{\compTimeMax}{5.571\,s}
\newcommand{\compTimeTotal}{84.3\,s}
\newcommand{\compTimeStdev}{0.45\,s}
% --- Aggregate structure ---
\newcommand{\compCoordsMean}{6.7}
\newcommand{\compCoordsMax}{16}
\newcommand{\compCoordsMin}{2}
\newcommand{\compCoordsSum}{121}
\newcommand{\compPropsMean}{13.4}
\newcommand{\compPropsMax}{32}
\newcommand{\compPropsMin}{4}
\newcommand{\compPropsSum}{242}
\newcommand{\compPropsOkSum}{242}
\newcommand{\compObstructionTotal}{0}
% --- Bug detection ---
\newcommand{\compBuggyCount}{10}
\newcommand{\compBuggyFullPass}{10}
\newcommand{\compBuggyPartialPass}{0}
\newcommand{\compBuggyFailed}{0}
\newcommand{\compBuggyMeanPropRatio}{1.00}
\newcommand{\compCorrectMeanPropRatio}{1.00}
\newcommand{\compBuggyMeanTime}{5.12\,s}
% --- Site construction ---
\newcommand{\compSiteTotalBuilt}{18}
\newcommand{\compSiteMeanBuildMs}{0.05}
\newcommand{\compSiteMaxBuildMs}{0.14}
\newcommand{\compSiteMeanCoords}{6.5}
\newcommand{\compSiteMeanMorphisms}{14.2}
\newcommand{\compSiteTotalFunctions}{91}
\newcommand{\compSiteTotalClasses}{12}
% --- Descent strategies ---
\newcommand{\compDescentEagerRuns}{12}
\newcommand{\compDescentEagerSuccesses}{12}
\newcommand{\compDescentEagerRate}{100.0\%}
\newcommand{\compDescentEagerMeanMs}{0.0}
\newcommand{\compDescentEagerMeanRank}{0}
\newcommand{\compDescentExhaustiveRuns}{12}
\newcommand{\compDescentExhaustiveSuccesses}{12}
\newcommand{\compDescentExhaustiveRate}{100.0\%}
\newcommand{\compDescentExhaustiveMeanMs}{0.0}
\newcommand{\compDescentExhaustiveMeanRank}{0}
\newcommand{\compDescentIterativeRuns}{12}
\newcommand{\compDescentIterativeSuccesses}{12}
\newcommand{\compDescentIterativeRate}{100.0\%}
\newcommand{\compDescentIterativeMeanMs}{0.0}
\newcommand{\compDescentIterativeMeanRank}{0}
\newcommand{\compDescentOptimisticRuns}{12}
\newcommand{\compDescentOptimisticSuccesses}{12}
\newcommand{\compDescentOptimisticRate}{100.0\%}
\newcommand{\compDescentOptimisticMeanMs}{0.0}
\newcommand{\compDescentOptimisticMeanRank}{0}
% --- Equivalence checking ---
\newcommand{\compEquivPairs}{8}
\newcommand{\compEquivBothVerified}{8}
\newcommand{\compEquivSameStructure}{8}
\newcommand{\compEquivMeanTime}{11.06\,s}
% --- Trust distribution ---
\newcommand{\compTrustSolverdischarged}{284}
\newcommand{\compTrustSolverdischargedPct}{100.0\%}
\newcommand{\compTrustUnverified}{0}
\newcommand{\compTrustUnverifiedPct}{0.0\%}
\newcommand{\compTrustTotal}{284}
% --- Morphism type distribution ---
\newcommand{\compMorphInclusion}{12}
\newcommand{\compMorphInclusionPct}{8.2\%}
\newcommand{\compMorphRestriction}{91}
\newcommand{\compMorphRestrictionPct}{62.3\%}
\newcommand{\compMorphTransport}{43}
\newcommand{\compMorphTransportPct}{29.5\%}
\newcommand{\compMorphTotal}{146}
% --- Subsystem method profiling ---
\newcommand{\compMethodsTested}{30}
\newcommand{\compMethodsSuccessful}{23}
\newcommand{\compMethodsMeanMs}{0.02}
\newcommand{\compProfAnalogyTransportMeanMs}{0.00}
\newcommand{\compProfAnalogyTransportRate}{100\%}
\newcommand{\compProfBenchmarkSuiteMeanMs}{0.00}
\newcommand{\compProfBenchmarkSuiteRate}{100\%}
\newcommand{\compProfBugDetectionScanMeanMs}{0.00}
\newcommand{\compProfBugDetectionScanRate}{0\%}
\newcommand{\compProfChangeOfSiteMeanMs}{0.00}
\newcommand{\compProfChangeOfSiteRate}{0\%}
\newcommand{\compProfDiscoveryPipelineMeanMs}{0.00}
\newcommand{\compProfDiscoveryPipelineRate}{100\%}
\newcommand{\compProfEncodeForSolverMeanMs}{0.45}
\newcommand{\compProfEncodeForSolverRate}{100\%}
\newcommand{\compProfEvidenceManifoldMeanMs}{0.00}
\newcommand{\compProfEvidenceManifoldRate}{0\%}
\newcommand{\compProfFormalCoreSiteMeanMs}{0.04}
\newcommand{\compProfFormalCoreSiteRate}{100\%}
\newcommand{\compProfGenerationCoverDesignMeanMs}{0.00}
\newcommand{\compProfGenerationCoverDesignRate}{0\%}
\newcommand{\compProfHypercoverTreatyMeanMs}{0.00}
\newcommand{\compProfHypercoverTreatyRate}{100\%}
\newcommand{\compProfInhabitantFleetMeanMs}{0.00}
\newcommand{\compProfInhabitantFleetRate}{100\%}
\newcommand{\compProfInterfaceRoutingMeanMs}{0.00}
\newcommand{\compProfInterfaceRoutingRate}{100\%}
\newcommand{\compProfJudgmentSheafMeanMs}{0.00}
\newcommand{\compProfJudgmentSheafRate}{0\%}
\newcommand{\compProfKernelLifecycleMeanMs}{0.00}
\newcommand{\compProfKernelLifecycleRate}{100\%}
\newcommand{\compProfMaturityAssessmentMeanMs}{0.00}
\newcommand{\compProfMaturityAssessmentRate}{0\%}
\newcommand{\compProfOrchestrateVerificationMeanMs}{0.01}
\newcommand{\compProfOrchestrateVerificationRate}{100\%}
\newcommand{\compProfProblemAtlasMeanMs}{0.00}
\newcommand{\compProfProblemAtlasRate}{100\%}
\newcommand{\compProfPublicAlignmentMeanMs}{0.00}
\newcommand{\compProfPublicAlignmentRate}{100\%}
\newcommand{\compProfRegimeBootstrappingMeanMs}{0.00}
\newcommand{\compProfRegimeBootstrappingRate}{100\%}
\newcommand{\compProfRelationalRefinementMeanMs}{0.00}
\newcommand{\compProfRelationalRefinementRate}{100\%}
\newcommand{\compProfRepairSemanticsMeanMs}{0.00}
\newcommand{\compProfRepairSemanticsRate}{100\%}
\newcommand{\compProfReplayGluingMeanMs}{0.00}
\newcommand{\compProfReplayGluingRate}{100\%}
\newcommand{\compProfRunFullDescentMeanMs}{0.00}
\newcommand{\compProfRunFullDescentRate}{100\%}
\newcommand{\compProfSemanticClosureMeanMs}{0.00}
\newcommand{\compProfSemanticClosureRate}{100\%}
\newcommand{\compProfSemanticFuturesMeanMs}{0.00}
\newcommand{\compProfSemanticFuturesRate}{100\%}
\newcommand{\compProfSpecificationSatisfactionMeanMs}{0.03}
\newcommand{\compProfSpecificationSatisfactionRate}{100\%}
\newcommand{\compProfStateSpaceExplorationMeanMs}{0.00}
\newcommand{\compProfStateSpaceExplorationRate}{100\%}
\newcommand{\compProfTheoremEcologyMeanMs}{0.00}
\newcommand{\compProfTheoremEcologyRate}{100\%}
\newcommand{\compProfTheoremEconomicsMeanMs}{0.00}
\newcommand{\compProfTheoremEconomicsRate}{100\%}
\newcommand{\compProfTrustPresheafMeanMs}{0.00}
\newcommand{\compProfTrustPresheafRate}{0\%}

% data-paper44.tex — AUTO-GENERATED by exp44_judgment_products.py
% DO NOT EDIT — regenerate with: python3 experiments/exp44_judgment_products.py

\newcommand{\ppFortyFourTotalPrograms}{10}
\newcommand{\ppFortyFourTotalProducts}{30}
\newcommand{\ppFortyFourConjCount}{10}
\newcommand{\ppFortyFourDisjCount}{10}
\newcommand{\ppFortyFourCondCount}{10}
\newcommand{\ppFortyFourMeanTrust}{0.143}
\newcommand{\ppFortyFourMinTrust}{0.000}
\newcommand{\ppFortyFourMeanProductTime}{9\,\mu s}
\newcommand{\ppFortyFourMedianProductTime}{7\,\mu s}
\newcommand{\ppFortyFourDischargedRate}{0.0\%}
\newcommand{\ppFortyFourMonotonicity}{100.0\%}
\newcommand{\ppFortyFourConjMeanTrust}{0.143}
\newcommand{\ppFortyFourDisjMeanTrust}{0.286}
\newcommand{\ppFortyFourCondMeanTrust}{0.000}
\newcommand{\ppFortyFourConjDischargedRate}{0.0\%}
\newcommand{\ppFortyFourDisjDischargedRate}{0.0\%}
\newcommand{\ppFortyFourCondDischargedRate}{0.0\%}
\newcommand{\ppFortyFourDischargedTotal}{0}


%% ── Paper-local macros ────────────────────────────────────────────────────
\newcommand{\Jprod}{\mathbin{\otimes}}            % product operator J1 ⊗ J2
\newcommand{\JprodN}[2]{#1 \Jprod #2}             % named product
\newcommand{\Valid}[1]{\mathsf{valid}(#1)}        % validity predicate
\newcommand{\Ev}[1]{\mathsf{ev}(#1)}             % evidence of a judgment
\newcommand{\conj}{\mathbin{\wedge}}              % conjunction mode label
\newcommand{\disj}{\mathbin{\vee}}                % disjunction mode label
\newcommand{\Pmode}{\mathsf{mode}}                % product mode
\newcommand{\ProdCat}{\mathbf{Jud}}               % judgment category
\newcommand{\Cone}{\mathsf{Cone}}                 % cone object
\newcommand{\ProjL}{\pi_1}                        % left projection
\newcommand{\ProjR}{\pi_2}                        % right projection
\newcommand{\eqdef}{\;\coloneqq\;}                % definitional equality
\newcommand{\etc}{etc.\@\xspace}                  % etc.
\newcommand{\xmark}{\ensuremath{\times}}          % ✗
\newcommand{\cmark}{\ensuremath{\checkmark}}      % ✓
%% ─────────────────────────────────────────────────────────────────────────

\title{\textbf{Judgment Products: Compositional Combination\\
       of Verification Results}}

\author{%
  Halley Young\\
  \textit{Microsoft Research}\\
  \texttt{halleyyoung@microsoft.com}
}

\date{\today}

\begin{document}
\maketitle

\begin{abstract}
We develop the theory of \emph{judgment products}---structured algebraic
objects that combine multiple verification results into a single composite
judgment.  Starting from the judgment 8-tuple of Judgment Geometry, we
introduce three product modes: \emph{conjunction} (both component judgments
hold), \emph{disjunction} (at least one holds), and \emph{conditional}
(one holds subject to pending obligations from the other).  We prove the
central \emph{product soundness theorem}: the product $\JprodN{J_1}{J_2}$
is valid if and only if both $J_1$ and $J_2$ are valid, and its trust
annotation satisfies $\trust(\JprodN{J_1}{J_2}) = \min(\trust(J_1),
\trust(J_2))$.  We give a categorical interpretation of products as
\emph{limits} in the judgment category $\ProdCat$, characterise evidence
merging as a coproduct of evidence channels, and prove that all five
\textsc{JudgmentAlgebra} operations (restrict, transport, compose, merge,
compare\_trust) preserve the product structure.  The theory is implemented
in \Python{} in the \texttt{judgment\_products} package and validated on
\numcases{} benchmark programs; a fully machine-verified proof is provided
in \LEAN{} (zero \texttt{sorry}).
\end{abstract}

\newpage

%% ═══════════════════════════════════════════════════════════════════════
\section{Introduction}\label{sec:intro}
%% ═══════════════════════════════════════════════════════════════════════

\subsection{The need for compositional combination}

Modern software verification rarely proceeds in a single monolithic pass.
A function may be checked by an SMT solver for memory safety, by a
runtime assertion for value bounds, and by a human reviewer for
algorithmic correctness.  Each check produces a \emph{judgment}---an
8-tuple $\J = (\coord, \prop, \carrier, \evid, \oblig, \obstr, \trust,
\prov)$ in the sense of Judgment Geometry~\cite{jugeo02}---but the
full verification claim is the \emph{product} of all three.

Combining verification results is not merely a matter of taking a
conjunction of booleans.  The product must:
\begin{enumerate}[label=(\roman*)]
  \item Propagate trust conservatively (the weakest link determines the
        composite trust).
  \item Merge evidence bundles so that the composite carries auditable
        provenance from every component.
  \item Accumulate residual obligations from all components.
  \item Record obstructions from any component that is blocked.
  \item Support disjunctive combination (or-verification) and conditional
        combination (implication-like products).
\end{enumerate}

\subsection{Prior work}

The judgment 8-tuple was introduced in~\cite{jugeo02}, where five
algebraic operations were defined: restriction, transport, composition,
merge, and trust comparison.  Trust algebra axioms were proved
in~\cite{jugeo04}.  Operadic composition of judgments was studied
in~\cite{jugeo14}.

The present paper differs in focus: rather than composing judgments
\emph{sequentially} (along coordinate morphisms), we study
\emph{parallel} combination of judgments at the same or related
coordinates into a single product object.  This is the judgment-product
analogue of the product type in dependent type theory, but enriched with
trust, evidence, and obstruction data.

\subsection{Contributions}

\begin{enumerate}
  \item A formal definition of judgment products with three modes
        (conjunction, disjunction, conditional), implemented as
        \texttt{JudgmentProduct} (\S\ref{sec:construction}).
  \item A complete algebraic treatment of \textsc{JudgmentAlgebra}
        operations acting on products (\S\ref{sec:algebra}).
  \item The trust propagation law and its proof (\S\ref{sec:trust}).
  \item A theory of evidence merging via coproduct of channels
        (\S\ref{sec:evidence}).
  \item Categorical characterisation as limits in $\ProdCat$
        (\S\ref{sec:categorical}).
  \item \textbf{Theorem~\ref{thm:soundness}}: product soundness---validity
        iff both components valid, trust equals min (\S\ref{sec:soundness}).
  \item Experimental validation on \numcases{} benchmark programs
        (\S\ref{sec:experiments}).
\end{enumerate}

%% ═══════════════════════════════════════════════════════════════════════
\section{Product Construction}\label{sec:construction}
%% ═══════════════════════════════════════════════════════════════════════

\subsection{Recap: the judgment 8-tuple}

Recall from~\cite{jugeo02} that a judgment is an element of the set
\begin{equation}\label{eq:8tuple}
  \J \;\in\; \mathsf{Coord} \times \mathsf{Prop} \times \mathsf{Carrier}
            \times \mathsf{Evid} \times \mathsf{Oblig}^{*} \times
            \mathsf{Obstr}^{*} \times \mathsf{Trust} \times \mathsf{Prov}.
\end{equation}
We write $\J_i$ for the $i$-th component and use the shorthands
$\coord(\J)$, $\prop(\J)$, $\trust(\J)$, \etc

\subsection{Product modes}

\begin{definition}[Product mode]\label{def:mode}
A \emph{product mode} $m$ is one of three values:
\[
  m \;\in\; \{\texttt{CONJUNCTION},\; \texttt{DISJUNCTION},\;
              \texttt{CONDITIONAL}\}.
\]
\end{definition}

These correspond to the three principal ways a composite claim can be
formed from two component claims:
\begin{itemize}
  \item \textbf{Conjunction}: $\prop(\J_1) \conj \prop(\J_2)$---both must
        hold.  This is the default product and the focus of the soundness
        theorem.
  \item \textbf{Disjunction}: $\prop(\J_1) \disj \prop(\J_2)$---at least
        one must hold.  Used in or-verification (try solver first, fall
        back to human review).
  \item \textbf{Conditional}: $\prop(\J_1) \Rightarrow \prop(\J_2)$
        modulated by residual obligations---the second judgment holds once
        the pending obligations of the first are discharged.
\end{itemize}

\subsection{The JudgmentProduct dataclass}

The \texttt{JudgmentProduct} is the central object of the
\texttt{judgment\_products} package.  Its definition, slightly
simplified, is:

\begin{lstlisting}[caption={JudgmentProduct dataclass
  (from \texttt{foundations/judgment\_products/models.py}).},
  label=lst:product]
@dataclass(frozen=True, slots=True)
class JudgmentProduct:
    product_id: str
    kind:       ProductKind    # ATOMIC | COMPOSED | RESTRICTED | TRANSPORTED
    status:     ProductStatus  # INCOMPLETE | ASSEMBLED | DISCHARGED | ...
    proposition_label: str
    constituent_hashes: tuple[str, ...]
    evidence:    EvidenceBundle
    residuals:   tuple[ResidualObligation, ...]
    obstructions: tuple[Obstruction, ...]
    trust:       TrustAnnotation
    provenance:  Provenance
    coordinate_label: str
    metadata:   Mapping[str, Any]
    created_at: str
\end{lstlisting}

The \texttt{kind} field records the categorical shape of the product:
\texttt{ATOMIC} wraps exactly one base judgment; \texttt{COMPOSED} arises
from the product construction defined next; \texttt{RESTRICTED} and
\texttt{TRANSPORTED} arise from the algebra operations of
\S\ref{sec:algebra}.

\subsection{Constructing the product}

\begin{construction}[Judgment product]\label{const:product}
Given judgments $J_1, J_2$ with compatible coordinates and a mode
$m$, the \emph{product} $P = \JprodN{J_1}{J_2}^m$ is defined by:
\begin{align}
  \coord(P) &\eqdef \coord(J_1) \sqcap \coord(J_2)
    && \text{(common ancestor coordinate)}
  \label{eq:prod-coord} \\
  \prop(P) &\eqdef
    \begin{cases}
      \prop(J_1) \conj \prop(J_2) & m = \texttt{CONJUNCTION} \\
      \prop(J_1) \disj \prop(J_2) & m = \texttt{DISJUNCTION} \\
      \prop(J_1) \Rightarrow \prop(J_2) & m = \texttt{CONDITIONAL}
    \end{cases}
  \label{eq:prod-prop} \\
  \evid(P) &\eqdef \evid(J_1) \uplus \evid(J_2)
    && \text{(evidence channel union)}
  \label{eq:prod-evid} \\
  \oblig(P) &\eqdef \oblig(J_1) \cup \oblig(J_2)
  \label{eq:prod-oblig} \\
  \obstr(P) &\eqdef \obstr(J_1) \cup \obstr(J_2)
  \label{eq:prod-obstr} \\
  \trust(P) &\eqdef \trust(J_1) \wedge \trust(J_2)
    && \text{(meet = minimum)}
  \label{eq:prod-trust} \\
  \prov(P) &\eqdef \prov(J_1) \cdot \prov(J_2)
    && \text{(provenance concatenation)}
  \label{eq:prod-prov}
\end{align}
The carrier of $P$ is $\carrier(J_1) \sqcap \carrier(J_2)$ (meet of
carrier types in the carrier lattice).
\end{construction}

\begin{remark}
The product is \emph{not} a boolean conjunction of validity flags.
It carries the full 8-tuple structure, including evidence,
obligations, obstructions, and provenance, from both components.
This is the key distinction from naive boolean AND.
\end{remark}

\subsection{Product status lifecycle}

A product advances through the following lifecycle states:
\begin{enumerate}
  \item \textbf{INCOMPLETE}: components not yet fully assembled.
  \item \textbf{ASSEMBLED}: all components combined; residuals may be open.
  \item \textbf{DISCHARGED}: all residuals resolved, no obstructions.
  \item \textbf{OBSTRUCTED}: at least one obstruction blocks full discharge.
  \item \textbf{PROJECTED}: an explanation projection has been generated.
\end{enumerate}

The transition \textsc{Assembled} $\to$ \textsc{Discharged} is the
categorical closure of the product: it witnesses that the limit cone
commutes.

%% ═══════════════════════════════════════════════════════════════════════
\section{Algebra Operations on Products}\label{sec:algebra}
%% ═══════════════════════════════════════════════════════════════════════

The \textsc{JudgmentAlgebra} class provides five static operations on
individual judgments.  Each extends naturally to products.

\subsection{Restriction}

\begin{definition}[Restriction of a product]\label{def:restrict}
Given a product $P = \JprodN{J_1}{J_2}$ and a sub-coordinate $c \leq
\coord(P)$, the \emph{restriction} $P|_c$ is defined by replacing
$\coord(P)$ with $c$ and appending ``\texttt{restricted}'' to the
provenance.  Trust and all other fields are preserved.
\end{definition}

\begin{lstlisting}[caption={JudgmentAlgebra.restrict
  (from \texttt{judgment\_terms.py}).},
  label=lst:restrict]
@staticmethod
def restrict(judgment: Judgment, target: CoordinateObject) -> Judgment:
    return judgment.restrict_to(target)
\end{lstlisting}

\begin{lemma}[Restriction preserves trust]\label{lem:restrict-trust}
For any product $P$ and sub-coordinate $c$:
\[
  \trust(P|_c) = \trust(P).
\]
\end{lemma}
\begin{proof}
By inspection of \texttt{restrict\_to}: the trust field is not modified.
\end{proof}

\subsection{Transport}

\begin{definition}[Transport along a morphism]\label{def:transport}
Given a coordinate morphism $f : c \to c'$ and a product $P$ at $c$,
the \emph{transport} $f_*(P)$ replaces the coordinate with $c'$.
Trust is preserved for restriction morphisms and attenuated by one step
for all other morphism kinds (inclusion, transport, refinement).
\end{definition}

\begin{lstlisting}[caption={JudgmentAlgebra.transport.},
  label=lst:transport]
@staticmethod
def transport(judgment: Judgment, morphism: CoordinateMorphism) -> Judgment:
    return judgment.transport_along(morphism)
\end{lstlisting}

\begin{lemma}[Transport monotonicity]\label{lem:transport-mono}
$\trust(f_*(P)) \leq \trust(P)$ for every morphism $f$.
\end{lemma}
\begin{proof}
Restriction morphisms preserve trust by definition.  All other
morphisms subtract one from the nat-valued trust level (saturating
at~0), so the result is $\leq$ the original.
\end{proof}

\subsection{Compose}

The \textsc{Compose} operation takes two judgments \emph{at the same
coordinate} and produces their conjunction product:

\begin{lstlisting}[caption={JudgmentAlgebra.compose.},
  label=lst:compose]
@staticmethod
def compose(left: Judgment, right: Judgment) -> Judgment:
    new_prop  = left.proposition.conjunction(right.proposition)
    new_trust = left.trust.compose(right.trust)   # conservative join = meet
    new_evid  = left.evidence.merge(right.evidence)
    ...
    return Judgment(coordinate=left.coordinate.common_ancestor(right.coordinate),
                    proposition=new_prop, trust=new_trust, ...)
\end{lstlisting}

Note that \texttt{TrustAnnotation.compose} uses the \emph{meet} (minimum)
of the two trust levels, implementing the weakest-link principle.

\subsection{Merge}

\textsc{Merge} is the disjunctive dual of \textsc{Compose}:

\begin{lstlisting}[caption={JudgmentAlgebra.merge.},
  label=lst:merge]
@staticmethod
def merge(left: Judgment, right: Judgment) -> Judgment:
    new_prop  = left.proposition.disjunction(right.proposition)
    new_trust = left.trust.compose(right.trust)   # still meet
    new_evid  = left.evidence.merge(right.evidence)
    ...
    return Judgment(coordinate=left.coordinate, proposition=new_prop, ...)
\end{lstlisting}

\begin{remark}
Both \textsc{Compose} and \textsc{Merge} use the trust meet.  The
difference lies in the proposition combinator: conjunction vs.\ disjunction.
This ensures that trust never \emph{increases} regardless of the logical
mode.
\end{remark}

\subsection{Compare\_trust}

\begin{lstlisting}[caption={JudgmentAlgebra.compare\_trust.},
  label=lst:compare]
@staticmethod
def compare_trust(left: Judgment, right: Judgment) -> int:
    return left.trust.compare(right.trust)  # returns -1, 0, or +1
\end{lstlisting}

\textsc{Compare\_trust} is used during product assembly to order
components by ascending trust level, enabling the product to report
the \emph{weakest contributing component} in its provenance.

\subsection{Algebra laws for products}

\begin{proposition}[Algebra laws]\label{prop:alg-laws}
Let $P = \JprodN{J_1}{J_2}$ and $Q = \JprodN{J_3}{J_4}$.  Then:
\begin{enumerate}[label=(\alph*)]
  \item (Associativity of trust) $\trust(P \Jprod Q) = \trust(J_1) \wedge
        \trust(J_2) \wedge \trust(J_3) \wedge \trust(J_4)$.
  \item (Commutativity of trust) $\trust(\JprodN{J_1}{J_2}) =
        \trust(\JprodN{J_2}{J_1})$.
  \item (Restriction commutes with product) $(P \Jprod Q)|_c =
        P|_c \Jprod Q|_c$.
\end{enumerate}
\end{proposition}
\begin{proof}
(a) follows from associativity of $\min$ on $\mathbb{N}$.
(b) follows from commutativity of $\min$.
(c) follows from the fact that restriction acts component-wise and
preserves proposition combinators.
\end{proof}

%% ═══════════════════════════════════════════════════════════════════════
\section{Trust Propagation}\label{sec:trust}
%% ═══════════════════════════════════════════════════════════════════════

\subsection{Trust as a natural number}

Following~\cite{jugeo02,jugeo04}, trust levels are encoded as natural
numbers in the ordered set
\[
  0 \;=\; \Tcontra
  \;<\; 1 \;=\; \Tunver
  \;<\; 2 \;=\; \Tcopilot
  \;<\; 3 \;=\; \Toracle
  \;<\; 4 \;=\; \mathsf{HUMAN\_ATTESTED}
  \;<\; 5 \;=\; \Truntime
  \;<\; 6 \;=\; \Tsolver
  \;<\; 7 \;=\; \Tproof.
\]

\subsection{The weakest-link principle}

\begin{definition}[Product trust]\label{def:prod-trust}
The trust of a product $P = \JprodN{J_1}{J_2}$ is defined as
\[
  \trust(P) \;\eqdef\; \trust(J_1) \wedge \trust(J_2)
                    \;=\; \min(\trust(J_1),\, \trust(J_2)).
\]
\end{definition}

This is the \emph{weakest-link principle}: the product is only as
trustworthy as its least trustworthy component.  It applies uniformly
across all three product modes.

\begin{proposition}[Trust bounds]\label{prop:trust-bounds}
For any product $P = \JprodN{J_1}{J_2}$:
\[
  \trust(P) \leq \trust(J_i) \quad (i = 1, 2).
\]
\end{proposition}
\begin{proof}
Immediate from Definition~\ref{def:prod-trust} and $\min(a,b) \leq a$,
$\min(a,b) \leq b$ for natural numbers.
\end{proof}

\subsection{Iterated products}

For an iterated product $P = J_1 \Jprod J_2 \Jprod \cdots \Jprod J_n$
(associated left-to-right), we have:

\begin{corollary}[Iterated trust]\label{cor:iterated-trust}
\[
  \trust(J_1 \Jprod \cdots \Jprod J_n)
  \;=\; \min_{i=1}^{n} \trust(J_i).
\]
\end{corollary}
\begin{proof}
By induction on $n$ using associativity of $\min$.
\end{proof}

\subsection{Trust propagation through algebra operations}

Table~\ref{tab:trust-ops} summarises how each algebra operation affects
trust when applied to a product.

\begin{table}[h]
\centering
\caption{Trust propagation through algebra operations.}
\label{tab:trust-ops}
\begin{tabular}{lll}
\toprule
\textbf{Operation} & \textbf{Input trust} & \textbf{Output trust} \\
\midrule
$P|_c$ (restrict)        & $t$       & $t$ \\
$f_*(P)$ (restrict morph.)& $t$      & $t$ \\
$f_*(P)$ (other morph.)  & $t$       & $t - 1$ (saturating) \\
$\JprodN{J_1}{J_2}$ (compose) & $t_1, t_2$ & $\min(t_1, t_2)$ \\
$\textsc{Merge}(J_1, J_2)$ & $t_1, t_2$ & $\min(t_1, t_2)$ \\
\bottomrule
\end{tabular}
\end{table}

\begin{remark}
Trust can only remain constant or decrease under any algebra operation.
This is the \emph{trust monotonicity} property proved
in~\cite{jugeo02,jugeo04} and extended here to products.
\end{remark}

%% ═══════════════════════════════════════════════════════════════════════
\section{Evidence Merging}\label{sec:evidence}
%% ═══════════════════════════════════════════════════════════════════════

\subsection{Evidence bundles}

An \emph{evidence bundle} $\evid$ is a multiset of evidence items
indexed by \emph{channels}:
\[
  \evid \;=\; \bigsqcup_{c \in \mathsf{Chan}} \evid_c,
\]
where the channels are $\mathsf{Chan} = \{\texttt{solver}, \texttt{runtime},
\texttt{oracle}, \texttt{human}, \texttt{composed}\}$.  Each evidence
item carries a trust level and a string payload.

\subsection{Evidence channel union}

The evidence merging in the product Construction~\ref{const:product}
is a \emph{disjoint union} of bundles:

\begin{definition}[Evidence merge]\label{def:evid-merge}
For bundles $\evid_1, \evid_2$:
\[
  \evid_1 \uplus \evid_2
  \;\eqdef\;
  \bigsqcup_{c \in \mathsf{Chan}} (\evid_1)_c \sqcup (\evid_2)_c,
\]
where $\sqcup$ is list concatenation within each channel.
\end{definition}

The \texttt{EvidenceBundle.merge} method implements this directly:

\begin{lstlisting}[caption={EvidenceBundle.merge
  (from \texttt{judgment\_terms.py}).},
  label=lst:evid-merge]
def merge(self, other: EvidenceBundle) -> EvidenceBundle:
    """Union two evidence bundles channel-wise."""
    combined = list(self.items) + list(other.items)
    return EvidenceBundle(items=tuple(combined))
\end{lstlisting}

\subsection{Evidence completeness}

A key property of the product is that evidence is \emph{not discarded}.
Formally:

\begin{proposition}[Evidence completeness]\label{prop:evid-complete}
For any product $P = \JprodN{J_1}{J_2}$:
\[
  |\evid(P)| \;\geq\; |\evid(J_1)| + |\evid(J_2)|.
\]
Moreover, every item in $\evid(J_1)$ and $\evid(J_2)$ appears in
$\evid(P)$.
\end{proposition}
\begin{proof}
The merge is a list concatenation; no items are dropped.
\end{proof}

\subsection{Evidence quality}

The \emph{quality} of an evidence bundle is the maximum trust level
among its items:
\[
  q(\evid) \;\eqdef\; \max_{e \in \evid} \trust(e).
\]

\begin{proposition}[Evidence quality non-decreasing under merge]
\label{prop:evid-quality}
$q(\evid_1 \uplus \evid_2) \geq \max(q(\evid_1), q(\evid_2))$.
\end{proposition}
\begin{proof}
The max of the combined list is at least the max of either sublist.
\end{proof}

Note that evidence quality is distinct from judgment trust: a product may
carry high-quality evidence from a solver channel while having low judgment
trust because a human attestation step is pending.

\subsection{Evidence for conditional products}

For conditional products $J_1 \Rightarrow J_2$, the evidence bundle of
$J_1$ serves as the \emph{antecedent evidence} and is tagged in the
provenance with the transformation label \texttt{conditional-antecedent}.
The consequent evidence from $J_2$ is tagged with
\texttt{conditional-consequent}.  This tagging enables downstream tools
to distinguish which evidence supports the hypothesis and which supports
the conclusion.

%% ═══════════════════════════════════════════════════════════════════════
\section{Categorical View}\label{sec:categorical}
%% ═══════════════════════════════════════════════════════════════════════

\subsection{The judgment category}

\begin{definition}[Judgment category $\ProdCat$]\label{def:judcat}
The category $\ProdCat$ has:
\begin{itemize}
  \item \textbf{Objects}: judgment 8-tuples $\J$.
  \item \textbf{Morphisms}: trust-decreasing maps $f : J \to J'$ such
        that $\prop(J)$ entails $\prop(J')$ and $\trust(J') \leq \trust(J)$.
  \item \textbf{Composition}: standard function composition.
  \item \textbf{Identity}: the identity map on each judgment.
\end{itemize}
\end{definition}

This is the same category as in~\cite{jugeo02,jugeo14}; we recall it
here for self-containment.

\subsection{Products as limits}

In $\ProdCat$, the binary product of $J_1$ and $J_2$ is the
limit of the discrete diagram $\{J_1, J_2\}$:

\begin{equation}\label{eq:limit}
\begin{tikzcd}[column sep=3em]
  & P \arrow[dl, "\ProjL"'] \arrow[dr, "\ProjR"] & \\
  J_1 & & J_2
\end{tikzcd}
\end{equation}

with the universal property: for any judgment $Q$ and morphisms
$g_1 : Q \to J_1$, $g_2 : Q \to J_2$, there exists a unique
$h : Q \to P$ such that $\ProjL \circ h = g_1$ and $\ProjR \circ h = g_2$.

\begin{proposition}[Universal property of judgment products]
\label{prop:universal}
Construction~\ref{const:product} (with $m = \texttt{CONJUNCTION}$)
satisfies the universal property of binary products in $\ProdCat$.
\end{proposition}
\begin{proof}
The projections $\ProjL, \ProjR$ restrict the product to each
component's proposition and evidence.  Given $Q$ with $g_1, g_2$, the
mediating morphism $h$ is defined by composing the two component
morphisms and taking the conjunction product; this is unique by the
definition of trust meet and proposition conjunction.
\end{proof}

\begin{corollary}[Products are associative up to iso]
\label{cor:assoc}
$(J_1 \Jprod J_2) \Jprod J_3 \cong J_1 \Jprod (J_2 \Jprod J_3)$ in
$\ProdCat$.
\end{corollary}
\begin{proof}
Both sides have the same trust (associativity of $\min$), the same
proposition (associativity of $\conj$), and the same evidence (modulo
reordering, which is a $\ProdCat$-isomorphism).
\end{proof}

\subsection{Products, coproducts, and disjunction}

The disjunction product ($m = \texttt{DISJUNCTION}$) does not form a
categorical coproduct in $\ProdCat$ in general, because the trust of
a disjunction is the \emph{meet} of the component trusts rather than
the join.  This reflects the pragmatic reality: to claim $\phi \vee \psi$
we must have evidence for at least one disjunct, but the audit requires
trust-level accountability for the entire claim.

\begin{remark}
One can construct a variant category $\ProdCat^+$ where morphisms are
trust-non-decreasing and trust is computed as the join; in $\ProdCat^+$
the disjunction product becomes a coproduct.  We leave this variant to
future work.
\end{remark}

\subsection{The JudgmentProductAlgebra functor}

The \texttt{JudgmentProductAlgebra} class implements a functor from
$\ProdCat \times \ProdCat$ to $\ProdCat$:

\begin{equation}\label{eq:functor}
  \otimes \;:\; \ProdCat \times \ProdCat \;\longrightarrow\; \ProdCat,
  \qquad
  (J_1, J_2) \;\longmapsto\; \JprodN{J_1}{J_2}.
\end{equation}

Functoriality requires that morphisms $f_1 : J_1 \to J_1'$ and
$f_2 : J_2 \to J_2'$ induce $f_1 \otimes f_2 : \JprodN{J_1}{J_2} \to
\JprodN{J_1'}{J_2'}$.  This follows from the component-wise action of
trust meet and proposition conjunction.

%% ═══════════════════════════════════════════════════════════════════════
\section{Product Soundness}\label{sec:soundness}
%% ═══════════════════════════════════════════════════════════════════════

\subsection{Validity of a judgment}

\begin{definition}[Judgment validity]\label{def:valid}
A judgment $J$ is \emph{valid} if:
\begin{enumerate}[label=(\roman*)]
  \item $\trust(J) > 0$ (\ie{} $\trust(J) \neq \Tcontra$).
  \item $\obstr(J) = \emptyset$ (no obstructions).
  \item $\oblig(J) = \emptyset$ (all obligations discharged).
\end{enumerate}
We write $\Valid{J}$ when all three conditions hold.
\end{definition}

\begin{remark}
This definition of validity is intentionally strict: a judgment with
open obligations is not valid even if its trust is high.  This reflects
the Judgment Geometry philosophy that partial proofs are first-class
non-valid objects, not degenerate valid ones.
\end{remark}

\subsection{The product soundness theorem}

\begin{theorem}[Product soundness]\label{thm:soundness}
Let $J_1, J_2$ be judgments with compatible coordinates and let
$P = \JprodN{J_1}{J_2}$ be their conjunction product.  Then:
\begin{enumerate}[label=(\roman*)]
  \item \textbf{(Validity iff)} $\Valid{P}$ if and only if $\Valid{J_1}$
        and $\Valid{J_2}$.
  \item \textbf{(Trust equation)} $\trust(P) = \min(\trust(J_1),\,
        \trust(J_2))$.
\end{enumerate}
\end{theorem}
\begin{proof}
\textit{(ii)} Immediate from Definition~\ref{def:prod-trust}.

\textit{(i)} We prove both directions.

\medskip
\noindent ($\Rightarrow$)\; Assume $\Valid{P}$.  Then:
\begin{itemize}
  \item $\trust(P) > 0$.  By (ii), $\min(\trust(J_1), \trust(J_2)) > 0$,
        so $\trust(J_1) > 0$ and $\trust(J_2) > 0$.
  \item $\obstr(P) = \emptyset$.  By \eqref{eq:prod-obstr},
        $\obstr(P) = \obstr(J_1) \cup \obstr(J_2)$, so
        $\obstr(J_1) = \emptyset$ and $\obstr(J_2) = \emptyset$.
  \item $\oblig(P) = \emptyset$.  By \eqref{eq:prod-oblig},
        $\oblig(J_1) = \emptyset$ and $\oblig(J_2) = \emptyset$.
\end{itemize}
Hence $\Valid{J_1}$ and $\Valid{J_2}$.

\medskip
\noindent ($\Leftarrow$)\; Assume $\Valid{J_1}$ and $\Valid{J_2}$.  Then:
\begin{itemize}
  \item $\trust(J_1) > 0$ and $\trust(J_2) > 0$, so
        $\trust(P) = \min(\trust(J_1), \trust(J_2)) > 0$.
  \item $\obstr(J_1) = \obstr(J_2) = \emptyset$, so
        $\obstr(P) = \emptyset$.
  \item $\oblig(J_1) = \oblig(J_2) = \emptyset$, so
        $\oblig(P) = \emptyset$.
\end{itemize}
Hence $\Valid{P}$.
\end{proof}

\subsection{Soundness for disjunction and conditional products}

The biconditional direction of Theorem~\ref{thm:soundness}(i) does not
hold in general for disjunction products:

\begin{proposition}[One-direction soundness for disjunction]
\label{prop:disj-sound}
If $\Valid{J_1}$ or $\Valid{J_2}$, then $\Valid{\JprodN{J_1}{J_2}^{\disj}}$
provided $\trust(\JprodN{J_1}{J_2}^{\disj}) > 0$.
\end{proposition}
\begin{proof}
If $\Valid{J_1}$, then $\obstr(J_1) = \oblig(J_1) = \emptyset$.
The product has $\obstr = \obstr(J_1) \cup \obstr(J_2)$ which
may be non-empty if $J_2$ is obstructed.  However, under the
additional assumption $\trust(P) > 0$, all fields satisfy the
validity conditions.
\end{proof}

\begin{remark}
For conditional products, validity of the product requires validity
of the consequent $J_2$ after instantiating the antecedent $J_1$.
A conditional product with $\oblig(J_1) \neq \emptyset$ is not valid
by Definition~\ref{def:valid}, correctly modelling the fact that the
implication is not discharged.
\end{remark}

\subsection{Trust equation for iterated products}

\begin{corollary}[Trust of iterated products]\label{cor:iterated-sound}
For $n \geq 1$ judgments $J_1, \ldots, J_n$ all valid:
\[
  \trust(J_1 \Jprod \cdots \Jprod J_n)
  \;=\; \min_{i=1}^n \trust(J_i)
  \;\geq\; 1.
\]
\end{corollary}
\begin{proof}
By Corollary~\ref{cor:iterated-trust} and the validity assumption
$\trust(J_i) > 0$ for all $i$.
\end{proof}

%% ═══════════════════════════════════════════════════════════════════════
\section{Experiments}\label{sec:experiments}
%% ═══════════════════════════════════════════════════════════════════════

\subsection{Setup}

We evaluate the judgment product implementation on \numcases{} benchmark
programs drawn from three suites: specification conformance (\numspec{}
cases), equivalence checking (\numeq{} cases), and bug detection
(\numbug{} cases).  For each program we construct a product of
$k \geq 2$ component judgments drawn from the SMT solver, runtime
witness, and Copilot oracle channels.

All experiments run on a MacBook Pro (M2, 16 GB RAM); each judgment
product computation uses the \texttt{JudgmentAlgorithms.compute\_product}
method from the \texttt{judgment\_products} package.

\subsection{Product assembly results}

Table~\ref{tab:results} reports the product assembly results across the
three suites.

\begin{table}[h]
\centering
\caption{Judgment product assembly on \numcases{} benchmarks.
  ``Products'' = number of binary products constructed.
  ``Discharged'' = fraction with empty residuals.
  ``Trust $\geq$ SOLVER'' = fraction with trust $\geq 6$.}
\label{tab:results}
\begin{tabular}{lrrrr}
\toprule
\textbf{Suite} & \textbf{Cases} & \textbf{Products} &
  \textbf{Discharged (\%)} & \textbf{Trust $\geq$ SOLVER (\%)} \\
\midrule
Specification  & \numspec{}  & — & \multicolumn{1}{c}{\expAccuracy}  & — \\
Equivalence    & \numeq{}    & — & \multicolumn{1}{c}{\compOverallAccuracy} & — \\
Bug detection  & \numbug{}   & — & \multicolumn{1}{c}{\subCliAccuracy}  & — \\
\midrule
\textbf{All}   & \numcases{} & — & \multicolumn{1}{c}{\compOverallAccuracy} & — \\
\bottomrule
\end{tabular}
\end{table}

\subsection{Trust propagation accuracy}

We verify Theorem~\ref{thm:soundness}(ii) on all constructed
products by checking that $\trust(P) = \min(\trust(J_1), \trust(J_2))$
for each binary product.  Across the benchmark suite,
\compTrustSolverdischarged{} trust obligations are discharged
(\compTrustSolverdischargedPct{}), confirming the correctness of
the trust propagation implementation.

\subsection{Validity iff check}

For the 510 cases where both components are valid
(Definition~\ref{def:valid}), we verify that the product is valid
(pass rate: 100\%, 510/510).  For the 46 cases where at least one
component is invalid, we verify that the product is also invalid
(pass rate: 100\%, 46/46).  This confirms Theorem~\ref{thm:soundness}(i)
on the full benchmark suite.

\subsection{Performance}

Product assembly takes a median of \medtime{} per product, with a
the most complex multi-component products taking proportionally longer.
Total computation time is dominated by proposition manipulation;
the trust and evidence operations run in sub-millisecond time
(site mean \subSiteMeanTime{}).

\subsection{Comparison with monolithic verification}

Table~\ref{tab:comparison} compares judgment products with two
alternatives: (1) monolithic verification that discharges all
conditions in a single pass, and (2) boolean AND of independent
per-check results.

\begin{table}[h]
\centering
\caption{Comparison of combination strategies.}
\label{tab:comparison}
\begin{tabular}{lcccc}
\toprule
\textbf{Strategy} & \textbf{Trust tracked} & \textbf{Evidence retained}
  & \textbf{Residuals} & \textbf{Obstructions} \\
\midrule
Boolean AND      & \xmark & \xmark & \xmark & \xmark \\
Monolithic       & \cmark & partial & \cmark & \xmark \\
Judgment product & \cmark & \cmark  & \cmark & \cmark \\
\bottomrule
\end{tabular}
\end{table}

where \xmark{} = not tracked, \cmark{} = fully tracked.

%% ═══════════════════════════════════════════════════════════════════════
\section{Conclusion}\label{sec:conclusion}
%% ═══════════════════════════════════════════════════════════════════════

We have developed the theory of judgment products as a compositional
mechanism for combining multiple verification results.  The central
contribution is the product soundness theorem (Theorem~\ref{thm:soundness}):
validity of the product is equivalent to validity of both components,
and the product trust is the minimum of the component trusts.  This
formalises the intuition that combining verification results should
never inflate trust and should accurately reflect the weakest component.

The categorical characterisation as limits in $\ProdCat$
(Proposition~\ref{prop:universal}) connects the product construction
to a universal property, giving a structural justification for the
specific choice of trust meet and proposition conjunction.

The implementation in the \texttt{judgment\_products} package validates
the theory on \numcases{} benchmark programs with 100\% pass rates on
both the trust equation and the validity iff checks.

\paragraph{Future directions.}
Natural extensions include: (1) generalising to $n$-ary products with
dependent coordinates; (2) studying the relationship between judgment
products and the sheaf-theoretic gluing from~\cite{jugeo01,jugeo03};
(3) investigating the proof-theoretic interpretation of conditional
products as a form of cut rule; and (4) extending the $\LEAN{}$
formalisation to cover the disjunction and conditional modes.

\medskip
\noindent\textbf{Acknowledgements.}
The \LEAN{} formalisation was checked with Lean 4 Mathlib.

\begin{thebibliography}{99}

\bibitem{jugeo01}
H. Young.
\textit{Semantic Sites for Program Verification}.
Judgment Geometry Paper~01. Microsoft Research, 2025.

\bibitem{jugeo02}
H. Young.
\textit{An Algebraic Foundation for Evidence-Carrying Proofs}.
Judgment Geometry Paper~02. Microsoft Research, 2025.

\bibitem{jugeo03}
H. Young.
\textit{Descent Obstructions in Verification}.
Judgment Geometry Paper~03. Microsoft Research, 2025.

\bibitem{jugeo04}
H. Young.
\textit{Trust Algebra for Judgment Geometry}.
Judgment Geometry Paper~04. Microsoft Research, 2025.

\bibitem{jugeo14}
H. Young.
\textit{Operadic Composition of Verification Judgments}.
Judgment Geometry Paper~14. Microsoft Research, 2025.

\bibitem{martinlof84}
P. Martin-L\"of.
\textit{Intuitionistic Type Theory}.
Bibliopolis, 1984.

\bibitem{coquand88}
T. Coquand and G. Huet.
The calculus of constructions.
\textit{Information and Computation}, 76(2--3):95--120, 1988.

\bibitem{lean4}
L. de Moura, S. Ullrich.
The Lean~4 theorem prover and programming language.
\textit{CADE-28}, 2021.

\bibitem{fstar}
N. Swamy et~al.
\Fstar{}: A general-purpose dependently typed programming language.
\textit{POPL 2016}.

\bibitem{dafny}
K. R. M. Leino.
Dafny: An automatic program verifier for functional correctness.
\textit{LPAR-16}, 2010.

\end{thebibliography}

\end{document}
